\documentclass[12pt,a4paper]{article}

% ── 中文支持 ──
\usepackage[UTF8]{ctex}
\setCJKmainfont{SimSun}[BoldFont=SimHei,ItalicFont=KaiTi]
\setCJKsansfont{SimHei}
\setCJKmonofont{FangSong}

% ── 页面 ──
\usepackage[top=2.5cm,bottom=2.5cm,left=2.5cm,right=2.5cm]{geometry}
\usepackage{setspace}
\onehalfspacing

% ── 数学 / 表格 / 代码 ──
\usepackage{amsmath,amssymb}
\usepackage{booktabs}
\usepackage{array}
\usepackage{tabularx}
\usepackage{longtable}
\usepackage{listings}
\usepackage{xcolor}
\usepackage{enumitem}
\usepackage{hyperref}
\hypersetup{colorlinks=true,linkcolor=black,citecolor=blue,urlcolor=blue}

% ── 代码样式 ──
\lstset{
  basicstyle=\small\ttfamily,
  backgroundcolor=\color{gray!8},
  frame=single,
  rulecolor=\color{gray!40},
  breaklines=true,
  numbers=left,
  numberstyle=\tiny\color{gray},
  keywordstyle=\color{blue!70},
  commentstyle=\color{green!50!black},
  stringstyle=\color{red!60},
  showstringspaces=false,
  tabsize=4,
}

% ── 标题层级 ──
\usepackage{titlesec}
\titleformat{\section}{\Large\bfseries\sffamily}{{\thesection}}{1em}{}
\titleformat{\subsection}{\large\bfseries\sffamily}{{\thesubsection}}{1em}{}
\titleformat{\subsubsection}{\normalsize\bfseries\sffamily}{{\thesubsubsection}}{1em}{}

\begin{document}

% ============================================================
%  封面
% ============================================================
\begin{titlepage}
\centering
\vspace*{4cm}
{\Huge\bfseries\sffamily 南京大学本科生课程设计报告 \par}
\vspace{1.5cm}
{\LARGE\bfseries 课程实验报告 \par}
\vspace{2cm}
{\Large
\begin{tabular}{rl}
\textbf{课程名称：} & 多智能体与强化学习 \\[6pt]
\textbf{选择题目：} & 峡谷漫步\;/\;重返秘境 \\[6pt]
\textbf{小组成员：} & 许政阳 231880081 \\
\end{tabular}
\par}
\vfill
\end{titlepage}

\tableofcontents
\newpage

% ============================================================
%  1  摘要
% ============================================================
\section{摘要}

本报告涵盖"多智能体与强化学习"课程的两次实验——"峡谷漫步"与"重返秘境"。两次实验均基于腾讯开悟（Kaiwu）强化学习平台，要求通过强化学习算法训练一个智能体，使其在包含障碍物的网格地图中学习移动策略，以最少的步数从起点导航到终点，并尽可能收集宝箱以获取更高分数。

实验一"峡谷漫步"（gorge\_walk）的核心目标是在 $64\times64$ 的网格地图中训练智能体完成基本的路径导航任务。本实验采用 Dueling DQN 算法，将 Q 值函数分解为状态价值函数 $V(s)$ 与优势函数 $A(s,a)$，通过共享特征层加双分支网络结构高效估计动作价值。训练过程中引入经验回放缓冲区和目标网络以增强稳定性，并结合课程学习、$\epsilon$ 指数衰减、多层次奖励塑造（包含距离引导、碰撞惩罚、振荡抑制和反停滞机制）等训练技巧。最终智能体以 64 步完成通关，获得 337 分。

实验二"重返秘境"（back\_to\_the\_realm）在更大的 $128\times128$ 地图上进行，动作空间扩展为 8 方向移动与技能使用的组合（16 维），同时引入宝箱收集、增益道具和闪烁技能等新机制。本实验采用 Target-DQN 算法并加入梯度裁剪策略，使用 CNN+MLP 混合网络结构处理包含空间信息的复合观测（404 维向量特征与 4 通道 $51\times51$ 局部地图），配合合法动作掩码、BFS 路径距离计算以及多维度奖励塑造。最终智能体在 237 步内完成通关，收集 4 个宝箱获得 400 分宝箱奖励，使用 1 次闪烁技能并获取 1 个增益道具，综合得分 702 分。

两次实验的对比展示了从简单导航到复杂多目标任务的渐进式学习过程，体现了算法选型、网络架构设计、特征工程和奖励塑造等关键要素在强化学习应用中的重要性。

% ============================================================
%  2  相关工作
% ============================================================
\section{相关工作}

深度强化学习（Deep Reinforcement Learning, DRL）将深度神经网络的表征学习能力与强化学习的决策优化框架相结合，近年来在游戏智能体、机器人控制、自动驾驶等领域取得了突破性进展。本节从值函数方法的演进、网格导航任务的研究以及奖励塑造技术三个维度梳理与本实验相关的工作。

\subsection{值函数方法的演进：从 DQN 到 Dueling DQN}

Mnih 等人\cite{mnih2015}提出的深度 Q 网络（DQN）首次成功地将深度神经网络应用于高维状态空间的值函数逼近。DQN 的两项核心创新——经验回放（Experience Replay）和目标网络（Target Network）——分别通过打破样本时间相关性和稳定训练目标，解决了深度网络训练 Q 函数时的不稳定问题。在此基础上，研究者提出了多种改进方案。

Van Hasselt 等人\cite{vanhasselt2016}提出 Double DQN，通过将动作选择与动作评估分离——使用在线网络选择最优动作、目标网络评估该动作的价值——有效缓解了标准 DQN 中的 Q 值过估计问题。Schaul 等人\cite{schaul2016}提出优先经验回放（Prioritized Experience Replay），根据 TD 误差对经验赋予不同的采样优先级，使得信息量更大的样本被更频繁地用于训练，显著提升了样本效率。

Wang 等人\cite{wang2016}提出的 Dueling DQN 是本实验一采用的核心算法。其创新在于将 Q 值函数显式分解为状态价值 $V(s)$ 和动作优势 $A(s,a)$ 两个流，即：
\begin{equation}
Q(s,a) = V(s) + A(s,a) - \frac{1}{|\mathcal{A}|}\sum_{a'}A(s,a')
\end{equation}
这种架构使网络能够独立学习哪些状态本身就有价值（无论采取何种动作），从而在动作对状态价值影响较小的场景下实现更高效的价值估计。Dueling DQN 在 Atari 游戏中的表现相比 DQN 有显著提升，尤其在动作空间较大但多数动作对结果影响相近的任务中表现突出。

\subsection{网格世界导航与路径规划}

网格世界（Grid World）是强化学习领域最经典的实验环境之一。从 Sutton 和 Barto\cite{sutton2018} 的教材中的简单迷宫到 MiniGrid（Chevalier-Boisvert 等, 2018）等复杂环境，网格导航任务一直是检验强化学习算法能力的重要基准。经典方法如 Q-learning、SARSA 等表格型方法可以解决小规模网格问题，但当状态空间增大（如本实验中的 $64\times64$ 和 $128\times128$ 网格）时，函数逼近方法成为必然选择。

在复杂网格导航中，局部观测和部分可观测性是常见挑战。研究者通常通过卷积神经网络（CNN）处理局部视野的空间信息，并结合循环网络（LSTM/GRU）或记忆机制来应对部分可观测问题。本实验二中采用的 CNN+MLP 混合架构正是遵循这一思路，使用 CNN 提取 $51\times51$ 局部视野中的障碍物、宝箱和终点等空间特征。此外，维护全局记忆图（memory map）记录访问历史也是一种常用的处理探索问题的手段，可视为一种简化的外部记忆机制。

\subsection{奖励塑造与课程学习}

奖励塑造（Reward Shaping）是将领域知识注入强化学习过程的重要手段。Ng 等人\cite{ng1999}从理论上证明了基于势函数的奖励塑造不会改变最优策略，为安全的奖励设计提供了理论基础。在实际应用中，距离引导奖励（鼓励靠近目标）、碰撞惩罚（抑制无效动作）、探索奖励（激励发现新区域）等是最常用的塑造策略。本实验中的多维度奖励设计综合运用了上述策略，尤其是实验二中按宝箱收集状态分阶段激活不同奖励项的方式，体现了层次化奖励设计的思想。

课程学习（Curriculum Learning, Bengio 等\cite{bengio2009}）通过由易到难地组织训练任务来加速收敛。在强化学习中，课程学习的典型实践包括逐步增大环境复杂度、分阶段引入任务目标等。本实验一中采用的"前 20\% 轮次关闭宝箱、先学通关再学收集"的策略即为课程学习的直接应用，先让智能体掌握基本导航能力，再增加宝箱收集的附属任务。

此外，梯度裁剪（Gradient Clipping）也是深度强化学习训练中的重要稳定化技术。Pascanu 等人\cite{pascanu2013}系统研究了循环网络中的梯度爆炸问题，提出通过限制梯度范数来稳定训练过程。本实验二中在 Target-DQN 的基础上引入 $\text{max\_norm}=10.0$ 的梯度裁剪，有效防止了 CNN 分支参数更新时可能出现的梯度爆炸。

% ============================================================
%  3  实验环境及实验设置
% ============================================================
\section{实验环境及实验设置}

\subsection{实验平台}

两次实验均在腾讯开悟（Kaiwu）强化学习平台上完成。开悟平台提供了基于浏览器的集成开发环境（IDE），支持在线代码编辑、训练任务管理、模型保存/加载、模型评估与回放等功能。实验采用 PyTorch 深度学习框架实现算法，通过平台提供的 \texttt{kaiwu\_agent} SDK 与环境进行交互。

% ────────────────────────────────────────────────────────────
%  3.2  实验一
% ────────────────────────────────────────────────────────────
\subsection{实验一：峡谷漫步}

\subsubsection{环境设置}

峡谷漫步环境（gorge\_walk）的基本设置如表~\ref{tab:hw1_env} 所示。

\begin{table}[htbp]
\centering
\caption{峡谷漫步环境参数}\label{tab:hw1_env}
\begin{tabular}{ll}
\toprule
\textbf{属性} & \textbf{描述} \\
\midrule
地图尺寸   & $64\times64$ 网格 \\
动作空间   & 4 个离散动作（上、下、左、右） \\
观测维度   & 250 维特征向量 \\
终止条件   & 到达终点（terminated）或超过最大步数（truncated） \\
\bottomrule
\end{tabular}
\end{table}

\subsubsection{算法设计}

本实验采用 \textbf{Dueling DQN} 算法。Dueling DQN 将 Q 值分解为状态价值 $V(s)$ 和优势函数 $A(s,a)$：
\begin{equation}
Q(s,a) = V(s) + A(s,a) - \text{mean}_{a'}A(s,a')
\end{equation}

\paragraph{网络架构}
网络包含共享特征层、Value 分支和 Advantage 分支，具体结构如表~\ref{tab:hw1_net} 所示。

\begin{table}[htbp]
\centering
\caption{Dueling DQN 网络结构}\label{tab:hw1_net}
\begin{tabular}{ll}
\toprule
\textbf{模块} & \textbf{结构} \\
\midrule
共享特征层       & Linear(250, 128) $\to$ ReLU \\
Value 分支      & Linear(128, 128) $\to$ ReLU $\to$ Linear(128, 1) \\
Advantage 分支  & Linear(128, 128) $\to$ ReLU $\to$ Linear(128, 4) \\
输出            & $Q = V + (A - \text{mean}(A))$ \\
\bottomrule
\end{tabular}
\end{table}

\paragraph{状态特征}
智能体的观测被处理为 250 维特征向量，包含以下信息：

\begin{itemize}[nosep]
\item \textbf{位置编码}（1 维）：将二维坐标编码为一维索引 $\text{pos\_x} \times 64 + \text{pos\_z}$
\item \textbf{行/列 One-hot}（$64+64=128$ 维）：当前行号和列号的 one-hot 编码
\item \textbf{距离特征}（$\sim$11 维）：到终点和各宝箱的离散化距离
\item \textbf{局部视野图}（$25\times3=75$ 维）：障碍物、宝箱、终点在局部视野中的分布
\item \textbf{记忆图}（25 维）：局部视野中已访问区域
\item \textbf{宝箱状态}（若干维）：各宝箱的收集状态
\end{itemize}

所有特征拼接后裁剪或零填充至 250 维。

\paragraph{经验回放与目标网络}
算法使用容量为 10000 的双端队列（deque）存储经验 $(s, a, r, s', \text{done})$，每次学习时随机采样 mini-batch（大小 64）进行训练。目标网络每隔 100 个学习步通过硬拷贝同步参数，TD 目标为：
\begin{equation}
y = r + (1 - \text{done}) \cdot \gamma \cdot \max_{a'} Q_{\text{target}}(s', a')
\end{equation}

\paragraph{奖励塑造}
设计了多层次的奖励塑造机制，如表~\ref{tab:hw1_reward} 所示。

\begin{table}[htbp]
\centering
\caption{峡谷漫步奖励设计}\label{tab:hw1_reward}
\begin{tabularx}{\textwidth}{lrX}
\toprule
\textbf{奖励项} & \textbf{值} & \textbf{说明} \\
\midrule
基础步数惩罚   & $-0.02$      & 鼓励智能体尽快完成任务 \\
靠近终点奖励   & $+0.08$      & 到终点距离减小时正奖励 \\
远离终点惩罚   & $-0.05$      & 到终点距离增大时负奖励 \\
靠近宝箱奖励   & $+0.02$      & 靠近最近有效宝箱时正奖励 \\
得分变化奖励   & $\Delta s/150$  & 环境分数变化缩放至 $[-1, 1]$ \\
通关奖励       & $+8.0$       & 成功到达终点 \\
超时惩罚       & $-8.0$       & 超过最大步数未通关 \\
原地惩罚       & $-0.15$      & 位置不变（碰撞障碍物） \\
两步往返惩罚   & $-0.12$      & 状态与上上步相同（振荡） \\
反向无效动作   & $-0.08$      & 执行反向动作且无收益 \\
停滞惩罚       & $-0.05$      & 连续 18 步无进展时触发随机探索 \\
\bottomrule
\end{tabularx}
\end{table}

\paragraph{训练策略}

\begin{enumerate}[nosep]
\item \textbf{Epsilon 衰减}：$\epsilon = \max(0.05,\; 1.0 \times 0.9995^{\text{episode}})$，从完全探索逐步过渡到利用。
\item \textbf{课程学习}：前 20\% 轮次关闭宝箱生成，先学导航再学收集。
\item \textbf{反停滞机制}：连续 18 步无进展时强制随机探索 4 步。
\item \textbf{早停}：滚动成功率 $> 97\%$ 且已完成至少 1000 episode 时提前结束。
\end{enumerate}

\subsubsection{超参数配置}

\begin{table}[htbp]
\centering
\caption{峡谷漫步超参数}\label{tab:hw1_hyper}
\begin{tabular}{llp{6cm}}
\toprule
\textbf{超参数} & \textbf{值} & \textbf{说明} \\
\midrule
学习率           & $10^{-3}$   & Adam 优化器 \\
折扣因子 $\gamma$  & 0.9       & 未来奖励折扣 \\
经验回放容量      & 10000     & 回放缓冲区最大容量 \\
批量大小          & 64        & 每次学习采样数 \\
目标网络更新      & 每 100 步  & 硬拷贝同步 \\
最大训练轮次      & 10000     & Episode 上限 \\
观测向量维度      & 250       & 输入特征维度 \\
\bottomrule
\end{tabular}
\end{table}

\subsubsection{实验结果与分析}

\begin{table}[htbp]
\centering
\caption{峡谷漫步实验结果}\label{tab:hw1_result}
\begin{tabular}{lr}
\toprule
\textbf{指标} & \textbf{值} \\
\midrule
得分         & 337 \\
步数         & 64 \\
宝箱收集数   & 0 \\
宝箱得分     & 0 \\
\bottomrule
\end{tabular}
\end{table}

智能体以 64 步完成通关，获得 337 分，表明 Dueling DQN 算法配合精心设计的奖励塑造能够有效学习峡谷漫步任务的导航策略。64 步的通关步数较为理想，说明智能体找到了一条高效路径。

宝箱收集数为 0，原因如下：(1) 课程学习前期关闭了宝箱，智能体优先学习了导航；(2) 通关奖励（$+8.0$）远大于宝箱奖励，策略倾向直奔终点；(3) 宝箱可能偏离最优路径，收集会增加步数。若要提升宝箱收集，可加大宝箱奖励权重或延长课程学习过渡期。

% ────────────────────────────────────────────────────────────
%  3.3  实验二
% ────────────────────────────────────────────────────────────
\subsection{实验二：重返秘境}

\subsubsection{环境设置}

重返秘境环境（back\_to\_the\_realm）的基本设置如表~\ref{tab:hw2_env} 所示。

\begin{table}[htbp]
\centering
\caption{重返秘境环境参数}\label{tab:hw2_env}
\begin{tabularx}{\textwidth}{lX}
\toprule
\textbf{属性} & \textbf{描述} \\
\midrule
地图尺寸   & $128\times128$ 网格（从 protobuf 格式的 map\_info 动态构建） \\
动作空间   & 16 个离散动作（8 方向 $\times$ 2 技能选项），带合法动作掩码 \\
观测维度   & 404 维向量 + 4 通道 $51\times51$ 局部特征图（共 4500 维） \\
附加机制   & 闪烁技能（Talent）、增益道具（Buff）、宝箱收集 \\
终止条件   & 到达终点（terminated）或超过最大步数（truncated） \\
\bottomrule
\end{tabularx}
\end{table}

\subsubsection{算法设计}

本实验采用 \textbf{Target-DQN + 梯度裁剪}。Target-DQN 使用独立的目标网络计算 TD 目标：
\begin{equation}
y = r + \gamma \cdot \max_{a'} Q_{\text{target}}(s', a') \cdot (1 - \text{done})
\end{equation}
目标网络每 500 步同步参数。反向传播后对梯度进行裁剪（$\text{clip\_grad\_norm\_}$ 上限 10.0），防止梯度爆炸。

\paragraph{网络架构}
模型采用 CNN+MLP 混合架构，具体结构如表~\ref{tab:hw2_net} 所示。

\begin{table}[htbp]
\centering
\caption{Target-DQN 混合网络结构}\label{tab:hw2_net}
\begin{tabularx}{\textwidth}{lX}
\toprule
\textbf{模块} & \textbf{结构} \\
\midrule
CNN 分支 & Conv2d(4,16,3,pad=2)$\to$BN$\to$ReLU$\to$MaxPool(2,2) $\to$ Conv2d(16,32,3,pad=2)$\to$BN$\to$ReLU$\to$MaxPool(2,2) $\to$ Conv2d(32,64,3,pad=2)$\to$BN$\to$ReLU$\to$MaxPool(2,2) $\to$ Flatten (4096维) \\
特征拼接 & CNN 输出 (4096) + 向量特征 (404) = 4500 维 \\
MLP 分支 & Linear(4500,256)$\to$ReLU$\to$Linear(256,128)$\to$ReLU$\to$Linear(128,16) \\
初始化   & Kaiming 正态初始化 \\
\bottomrule
\end{tabularx}
\end{table}

CNN 分支从 4 通道局部特征图（障碍物图、终点图、宝箱图、记忆图）中提取空间结构信息，MLP 将 CNN 输出与向量特征拼接后映射为 16 个动作的 Q 值。

\paragraph{状态特征}
观测特征分为向量部分和地图部分：

\textbf{向量特征（404 维）：}
\begin{itemize}[nosep]
\item 归一化位置（2 维）：英雄坐标归一化到 $[0,1]$
\item 位置 One-hot（256 维）：$128+128$ 维行列编码
\item 目标引导特征（9 维）：8 方向独热 + 归一化网格距离（当存在有效宝箱时动态指向最近宝箱，否则指向终点）
\item 宝箱相对特征（135 维）：15 个宝箱各 9 维
\item 增益/技能状态（2 维）：Buff 和 Talent 可用性
\end{itemize}

\textbf{地图特征（$4\times51\times51$）：}
\begin{itemize}[nosep]
\item 障碍物图：局部视野中的可通行/不可通行区域
\item 终点图：终点在局部视野中的位置标记
\item 宝箱图：当前存活宝箱的位置标记
\item 记忆图：已访问区域热度图
\end{itemize}

\texttt{Preprocessor} 类负责特征预处理：从 protobuf 原始观测中解析信息，维护 $128\times128$ 全局记忆图，使用 BFS 计算到各目标的网格距离，维护最近 100 步位置记录。

\paragraph{奖励塑造}
重返秘境的奖励函数包含多维度加权信号，如表~\ref{tab:hw2_reward} 所示。

\begin{table}[htbp]
\centering
\caption{重返秘境奖励设计}\label{tab:hw2_reward}
\begin{tabularx}{\textwidth}{lrX}
\toprule
\textbf{奖励项} & \textbf{权重} & \textbf{说明} \\
\midrule
终点距离奖励   & 0.5    & 收集 $\geq1$ 个宝箱后，靠近/远离终点给予 $\pm1$ \\
通关奖励       & 0.5    & 到达终点时，值与已收集宝箱数成正比 \\
宝箱距离奖励   & 0.5    & 仍有宝箱未收集时，靠近/远离最近宝箱给予 $\pm1$ \\
宝箱收集奖励   & 1.0    & 每成功收集一个宝箱 \\
步数惩罚       & $-0.0005$ & 微小的每步惩罚 \\
碰撞惩罚       & $-1.0$   & 两帧间移动距离过小（$\leq$500，即碰撞或原地踏步） \\
记忆惩罚       & $-0.005$ & 重复访问已走过区域 \\
探索奖励       & 0.1    & 首次到达新位置给正奖励；重复到达给递增惩罚 \\
超时惩罚       & $-3.0$   & 超过最大步数未通关（DIY 改进） \\
\bottomrule
\end{tabularx}
\end{table}

核心思路是优先收集宝箱（权重 1.0 最高），收集后再导航至终点（收集 $\geq1$ 个宝箱后激活终点奖励）。超时惩罚（$-3.0$）是相对于基线的改进，有效教导智能体避免浪费时间。

\subsubsection{超参数配置}

\begin{table}[htbp]
\centering
\caption{重返秘境超参数}\label{tab:hw2_hyper}
\begin{tabular}{llp{6cm}}
\toprule
\textbf{超参数} & \textbf{值} & \textbf{说明} \\
\midrule
学习率              & $10^{-4}$    & Adam，较小以配合更大网络 \\
折扣因子 $\gamma$     & 0.95       & 更重视长期回报 \\
$\epsilon$ 初始值    & 0.1        & 探索率 \\
$\epsilon$ 衰减     & 线性 / 30万步  & 约 30 万步后接近下界 \\
目标网络更新         & 每 500 步    & 较大间隔以稳定训练 \\
梯度裁剪上限         & 10.0       & 梯度最大范数 \\
局部视野             & $51\times51$ & VIEW\_SIZE = 25 \\
样本截断长度         & 256        & 每批次样本长度 \\
\bottomrule
\end{tabular}
\end{table}

\subsubsection{实验结果与分析}

\begin{table}[htbp]
\centering
\caption{重返秘境实验结果}\label{tab:hw2_result}
\begin{tabular}{lr}
\toprule
\textbf{指标} & \textbf{值} \\
\midrule
得分           & 702 \\
步数           & 237 \\
宝箱收集数     & 4 \\
宝箱得分       & 400 \\
技能使用次数   & 1 \\
增益获取数     & 1 \\
\bottomrule
\end{tabular}
\end{table}

智能体综合得分 702 分，各维度分析如下：

\begin{enumerate}[nosep]
\item \textbf{宝箱收集}：成功收集 4 个宝箱（400 分），验证了以宝箱收集为优先的奖励设计。"终点奖励在收集 $\geq1$ 个宝箱后才激活"的设计成功引导了先收集后导航的策略。
\item \textbf{导航效率}：237 步完成通关，在 $128\times128$ 大地图上绕道收集 4 个宝箱后仍能到达终点。
\item \textbf{技能与增益}：使用 1 次闪烁技能、获取 1 个增益道具，表明一定的技能利用能力。
\item \textbf{综合}：CNN 空间特征提取 + BFS 距离引导 + 分阶段奖励激活是取得良好表现的关键。
\end{enumerate}

% ────────────────────────────────────────────────────────────
%  3.4  对比
% ────────────────────────────────────────────────────────────
\subsection{两次实验对比}

\begin{table}[htbp]
\centering
\caption{两次实验对比}\label{tab:compare}
\begin{tabularx}{\textwidth}{lXX}
\toprule
\textbf{对比维度} & \textbf{实验一：峡谷漫步} & \textbf{实验二：重返秘境} \\
\midrule
算法     & Dueling DQN + 经验回放         & Target-DQN + 梯度裁剪 \\
网络     & MLP (250$\to$128$\to$Q)        & CNN+MLP (图+向量$\to$Q) \\
动作空间 & 4（上下左右）                   & 16（8 方向$\times$2 技能） \\
地图     & $64\times64$                    & $128\times128$ \\
得分     & 337                             & 702 \\
步数     & 64                              & 237 \\
宝箱     & 0 个 / 0 分                     & 4 个 / 400 分 \\
\bottomrule
\end{tabularx}
\end{table}

% ============================================================
%  4  总结和展望
% ============================================================
\section{总结和展望}

\subsection{总结}

通过两次实验的完整实践，我们对深度强化学习在导航任务中的应用获得了系统性的认识。实验一中，Dueling DQN 配合经验回放和目标网络成功解决了 $64\times64$ 网格中的路径规划问题，课程学习和多层次奖励塑造有效加速了训练收敛。实验二中，面对更大的状态空间、更复杂的动作空间和多目标任务，Target-DQN 结合 CNN 空间特征提取、合法动作掩码和梯度裁剪等技术实现了较好的综合表现。

两次实验的核心经验可归纳为：\textbf{第一}，算法选型应匹配问题复杂度——简单任务适合轻量级 MLP，复杂任务需要 CNN 等更强的特征提取能力；\textbf{第二}，奖励塑造是引导智能体行为的关键杠杆，奖励权重直接决定策略偏好（实验一中通关奖励主导致忽视宝箱，实验二中宝箱奖励优先则成功引导收集行为）；\textbf{第三}，训练稳定性技术（经验回放、目标网络、梯度裁剪）不可或缺；\textbf{第四}，特征工程仍然重要——BFS 距离、记忆图、位置编码等为智能体提供了充分的决策依据。

\subsection{不足}

当前实验存在一些不足：实验一中智能体未能收集宝箱，反映出奖励设计在多目标权衡上的欠缺；实验二虽收集了 4 个宝箱，但仍有提升空间。此外，两次实验均未引入 Double DQN 的动作选择-评估分离机制，可能存在 Q 值过估计；实验二未使用经验回放，样本利用效率有待提升；CNN 分支较为简单，对大范围空间信息的捕捉能力有限。

\subsection{展望}

未来改进方向包括：(1) 引入 Double DQN 减少过估计、Prioritized Experience Replay 提升样本效率、Noisy Network 实现更自适应的探索；(2) 探索注意力机制使智能体动态关注关键区域，或引入 LSTM/GRU 处理时序依赖；(3) 采用更精细的课程学习方案或分层强化学习，将宝箱收集和终点导航分为高层子目标；(4) 探索逆向强化学习从专家行为中自动学习奖励函数，避免手工调参的繁琐过程。

% ============================================================
%  参考文献
% ============================================================
\begin{thebibliography}{9}

\bibitem{mnih2015}
Mnih V, Kavukcuoglu K, Silver D, et al.
Human-level control through deep reinforcement learning[J].
\textit{Nature}, 2015, 518(7540): 529--533.

\bibitem{vanhasselt2016}
Van Hasselt H, Guez A, Silver D.
Deep reinforcement learning with double Q-learning[C].
\textit{AAAI}, 2016.

\bibitem{wang2016}
Wang Z, Schaul T, Hessel M, et al.
Dueling network architectures for deep reinforcement learning[C].
\textit{ICML}, 2016.

\bibitem{schaul2016}
Schaul T, Quan J, Antonoglou I, et al.
Prioritized experience replay[C].
\textit{ICLR}, 2016.

\bibitem{ng1999}
Ng A Y, Harada D, Russell S.
Policy invariance under reward transformations: Theory and application to reward shaping[C].
\textit{ICML}, 1999.

\bibitem{bengio2009}
Bengio Y, Louradour J, Collobert R, et al.
Curriculum learning[C].
\textit{ICML}, 2009.

\bibitem{sutton2018}
Sutton R S, Barto A G.
Reinforcement learning: An introduction[M].
MIT Press, 2018.

\bibitem{pascanu2013}
Pascanu R, Mikolov T, Bengio Y.
On the difficulty of training recurrent neural networks[C].
\textit{ICML}, 2013.

\end{thebibliography}

\end{document}
